\documentclass[11pt]{article}
\usepackage[margin=2.3cm]{geometry}
\usepackage{amsmath,amssymb,graphicx,listings,booktabs}
\newcommand{\fig}[2]{\begin{figure}[ht]\centering\includegraphics[width=\linewidth]{../figures/#1}\caption{#2}\end{figure}}
\title{Project 8: Inverse Problems with Multiple Valid Solutions}
\date{November 21, 2026}
\author{MAC 312}
\begin{document}
\maketitle

\section*{Setup}
$t\sim\mathrm{Unif}[-a,a]$ with $a=2$, $x=t^2+\epsilon$, $\epsilon\sim\mathcal N(0,\sigma^2)$, $\sigma=0.2$ unless stated. Train/validation/test sizes are 5000/1000/2000, with data seed 0 in every experiment; the seed in the experiment matrix controls only initialization and batch order. All three models share one body, $1\to64\to64$ with $\tanh$, trained with Adam (lr $10^{-3}$, batches of 256, up to 3000 epochs, early stopping on validation loss with patience 200, gradient clipping at 5). The single-Gaussian network is the MDN with $M=1$. Test data is never used for model selection.

\section{Part 1: exact reference}
\paragraph{Derivation.} With prior $p(t)=1/(2a)$ on $[-a,a]$, Bayes' rule gives
\[
p(t\mid x)=\frac{p(x\mid t)p(t)}{\int p(x\mid t')p(t')\,dt'}\propto \exp\!\Big(-\frac{(x-t^2)^2}{2\sigma^2}\Big),\quad |t|\le a .
\]
For an asymmetric prior with $P(t<0)=q$ the prior factor is $q/a$ for $t<0$ and $(1-q)/a$ for $t\ge0$. The constant is computed by the trapezoid rule on 2001 grid points (\texttt{scipy.integrate.trapezoid}, the same rule as \texttt{np.trapz}, which is removed in NumPy 2.x). Normalization is done in log space (subtract the per-row maximum before exponentiating), so that $x$ far outside the training range ($x=10$ or $x=-1$) gives a finite density rather than $0/0$.

\paragraph{Regions.}
\begin{itemize}
\item $x\le0$: one blob centred at $t=0$ (for $x<0$ the best $t^2$ is $0$, the least-wrong value).
\item $x$ clearly $>0$: two peaks at $\pm\sqrt{x}$, each of width about $\sigma/(2\sqrt x)$, with equal mass under the symmetric prior. For $x$ near $0$ (e.g.\ $x=0.1$) the two peaks overlap and merge into a flat top, so one versus two peaks is hard to tell.
\item $x$ near $a^2=4$: the peaks at $\pm\sqrt x$ are pushed against $\pm a$ and truncated by the support.
\item $x\gg a^2$ (e.g.\ $x=4.5$): all mass piles up at $t=\pm a$; the density has mass at the boundary only, and any learned model is extrapolating beyond the training data.
\end{itemize}
\fig{part1_reference.png}{Reference $p(t\mid x)$ for $x\in\{-0.2,0.1,1,3,4.5\}$; dotted lines mark $\pm\sqrt x$.}

\section{Part 2: the MSE network}
Minimizing $\mathbb E[(t-f(x))^2]$ over functions $f$: for each $x$, $\mathbb E[(t-f)^2\mid x]=\mathrm{Var}(t\mid x)+(f-\mathbb E[t\mid x])^2$, minimized by $f(x)=\mathbb E[t\mid x]$. Under the symmetric prior $p(t\mid x)$ is symmetric about $0$, so $\mathbb E[t\mid x]=0$ for every $x$: the network answers $0$, in the valley between the two branches, where there is no probability.
\fig{part2_mse.png}{Training data, the MSE network (red) and the reference $\mathbb E[t\mid x]$ (dashed).}

\section{Part 3: the mixture density network}
The head outputs $3M$ numbers per $x$: logits for $\pi_j$, means $\mu_j$ and raw variances. Weights use softmax, variances use $v_j=\mathrm{softplus}(\text{raw}_j)+10^{-4}$. The loss is the negative log-likelihood,
\[
\mathcal L=-\sum_{n}\log\sum_{j=1}^{M}\pi_j(x_n)\mathcal N(t_n\mid\mu_j(x_n),v_j(x_n)),
\]
which we divide by $N$ so that its value does not depend on batch size. It is computed as
$\mathcal L=-\sum_n\operatorname{logsumexp}_j\big[\log\pi_j-\tfrac12\log(2\pi v_j)-\tfrac{(t_n-\mu_j)^2}{2v_j}\big]$ with $\log\pi=\texttt{log\_softmax}(\text{logits})$. Working in log space and using \texttt{logsumexp} means a point far from every component gives a large but finite loss instead of $\log 0=-\infty$; the softplus keeps $v>0$ and the floor stops a component collapsing onto a single point (which would send the loss to $-\infty$). Means are initialized spread over $[-a,a]$ to discourage collapse onto one branch. Derivatives come from PyTorch autograd.

\section{Part 4: comparison}
Log-likelihood and density shape are the primary metrics. Test MSE is shown only to show that it misleads.
\begin{table}[ht]\centering\begin{tabular}{lrr}
\toprule
Model & Test mean log-lik. & Test MSE \\
\midrule
Exact reference & 0.184 & 1.296 \\
MSE network & --- & 1.300 \\
Single Gaussian & -1.235 & 1.299 \\
MDN & 0.163 & 1.298 \\
\bottomrule
\end{tabular}
\caption{Seed-0 models on the test set; the exact row uses the quadrature reference.}\end{table}
\begin{table}[ht]\centering\small\resizebox{\linewidth}{!}{\begin{tabular}{lrrrrrr}
\toprule
kind & log-lik. & MSE & branch err. & resid. std & rare branch & diverged \\
\midrule
mse & --- & 1.299 $\pm$ 0.001 & --- & --- & --- & 0/5 \\
gauss & -1.235 $\pm$ 0.000 & 1.298 $\pm$ 0.001 & 0.023 $\pm$ 0.004 & 2.472 $\pm$ 0.013 & 5/5 & 0/5 \\
mdn & 0.164 $\pm$ 0.003 & 1.299 $\pm$ 0.003 & 0.040 $\pm$ 0.010 & 0.206 $\pm$ 0.003 & 5/5 & 0/5 \\
\bottomrule
\end{tabular}
}\caption{Baseline comparison over 5 seeds (mean $\pm$ std).}\end{table}
\fig{part4_overlays.png}{Reference (black), single Gaussian (blue dashed) and MDN (red).}
\fig{part4_branch_prob.png}{$P(t>0\mid x)$, reference vs MDN.}
\fig{part4_residuals.png}{Forward residuals $x-t_s^2$ for sampled $t_s$, against $\mathcal N(0,\sigma^2)$.}
The residuals are compared with $\mathcal N(0,\sigma^2)$ and not with zero: if a sampled $t_s$ is correct, then $x-t_s^2$ should still contain the measurement noise $\epsilon$, so a spread of exactly zero would mean overfitting the noise.

\section{Part 5: stress tests}
Each row below is the mean $\pm$ standard deviation over 5 seeds. ``Rare branch'' is the number of runs in which the model gives $P(t<0\mid x=2)$ at least half the reference value; ``diverged'' counts runs that produced a non-finite loss.
\paragraph{Noise level.}\begin{table}[ht]\centering\small\resizebox{\linewidth}{!}{\begin{tabular}{lrrrrrr}
\toprule
sigma & log-lik. & MSE & branch err. & resid. std & rare branch & diverged \\
\midrule
0.05 & 1.384 $\pm$ 0.004 & 1.296 $\pm$ 0.001 & 0.032 $\pm$ 0.003 & 0.055 $\pm$ 0.001 & 5/5 & 0/5 \\
0.2 & 0.164 $\pm$ 0.003 & 1.299 $\pm$ 0.003 & 0.040 $\pm$ 0.010 & 0.206 $\pm$ 0.003 & 5/5 & 0/5 \\
0.5 & -0.551 $\pm$ 0.002 & 1.298 $\pm$ 0.003 & 0.022 $\pm$ 0.007 & 0.508 $\pm$ 0.002 & 5/5 & 0/5 \\
\bottomrule
\end{tabular}
}\end{table}
\fig{part5_noise.png}{MDN vs reference at $x=1$ for each $\sigma$.}
\paragraph{Asymmetric prior.}\begin{table}[ht]\centering\small\resizebox{\linewidth}{!}{\begin{tabular}{lrrrrrr}
\toprule
p neg & log-lik. & MSE & branch err. & resid. std & rare branch & diverged \\
\midrule
0.5 & 0.164 $\pm$ 0.003 & 1.299 $\pm$ 0.003 & 0.040 $\pm$ 0.010 & 0.206 $\pm$ 0.003 & 5/5 & 0/5 \\
0.2 & 0.377 $\pm$ 0.001 & 0.825 $\pm$ 0.002 & 0.037 $\pm$ 0.017 & 0.209 $\pm$ 0.002 & 5/5 & 0/5 \\
0.05 & 0.693 $\pm$ 0.002 & 0.290 $\pm$ 0.000 & 0.015 $\pm$ 0.005 & 0.209 $\pm$ 0.004 & 5/5 & 0/5 \\
\bottomrule
\end{tabular}
}\end{table}
\fig{part5_prior.png}{Branch probability and density at $x=1$ for each $P(t<0)$ (solid: reference, dashed: MDN).}
\paragraph{Number of components.}\begin{table}[ht]\centering\small\resizebox{\linewidth}{!}{\begin{tabular}{lrrrrrr}
\toprule
M & log-lik. & MSE & branch err. & resid. std & rare branch & diverged \\
\midrule
1 & -1.235 $\pm$ 0.000 & 1.298 $\pm$ 0.001 & 0.023 $\pm$ 0.004 & 2.472 $\pm$ 0.013 & 5/5 & 0/5 \\
2 & 0.152 $\pm$ 0.001 & 1.295 $\pm$ 0.001 & 0.027 $\pm$ 0.005 & 0.209 $\pm$ 0.002 & 5/5 & 0/5 \\
3 & 0.164 $\pm$ 0.003 & 1.299 $\pm$ 0.003 & 0.040 $\pm$ 0.010 & 0.206 $\pm$ 0.003 & 5/5 & 0/5 \\
5 & 0.164 $\pm$ 0.002 & 1.298 $\pm$ 0.002 & 0.029 $\pm$ 0.010 & 0.196 $\pm$ 0.004 & 5/5 & 0/5 \\
\bottomrule
\end{tabular}
}\end{table}
\fig{part5_components.png}{Density at $x=1,3$ for each $M$.}
\paragraph{Variance floor.}\begin{table}[ht]\centering\small\resizebox{\linewidth}{!}{\begin{tabular}{lrrrrrr}
\toprule
var floor & log-lik. & MSE & branch err. & resid. std & rare branch & diverged \\
\midrule
0.0 & 0.161 $\pm$ 0.004 & 1.302 $\pm$ 0.008 & 0.045 $\pm$ 0.013 & 0.208 $\pm$ 0.003 & 5/5 & 0/5 \\
0.0001 & 0.164 $\pm$ 0.003 & 1.299 $\pm$ 0.003 & 0.040 $\pm$ 0.010 & 0.206 $\pm$ 0.003 & 5/5 & 0/5 \\
0.01 & 0.091 $\pm$ 0.001 & 1.299 $\pm$ 0.002 & 0.042 $\pm$ 0.006 & 0.269 $\pm$ 0.001 & 5/5 & 0/5 \\
\bottomrule
\end{tabular}
}\end{table}
\fig{part5_floor.png}{Validation NLL and density at $x=1$ for each floor.}
\paragraph{Initialization.} Every row above already averages over 5 initialization seeds; the $\pm$ values are the initialization spread.

\paragraph{Failures.} The numbers of diverged runs and of runs that lost the rare branch are in the last two columns of the tables above, and are the only counts we claim.

\paragraph{Why MSE cannot assess a multimodal answer.} Squared error scores only the conditional mean $\mathbb E[t\mid x]$. A model with one peak at $0$, a model with two peaks at $\pm\sqrt x$ and the exact posterior all have the same mean, so under the symmetric prior they all get the same test MSE, as the test-MSE column of the first baseline table and the tables above shows, while their log-likelihoods differ greatly. A single number about the mean cannot distinguish a unimodal from a bimodal distribution.

\clearpage
\section*{Appendix: code listings}
\lstset{basicstyle=\ttfamily\scriptsize, breaklines=true, numbers=left, numberstyle=\tiny, frame=single}
\subsection*{data.py}
\lstinputlisting[language=Python]{../data.py}
\subsection*{reference.py}
\lstinputlisting[language=Python]{../reference.py}
\subsection*{models.py}
\lstinputlisting[language=Python]{../models.py}
\subsection*{train.py}
\lstinputlisting[language=Python]{../train.py}
\subsection*{evaluate.py}
\lstinputlisting[language=Python]{../evaluate.py}

\end{document}
